\documentclass[11pt,a4paper]{article}

\usepackage[utf8]{inputenc}
\usepackage[T1]{fontenc}
\usepackage[french]{babel}
\usepackage{geometry}
\usepackage{hyperref}

\geometry{margin=2.5cm}

\title{P6.1 - Gestion d'utilisateurs NFS/NIS}
\author{Yann RENVOISE \and Valentin MENON}
\date{\today}

\begin{document}

\maketitle
\tableofcontents

\section{Introduction}

Ce rapport presente un projet d'administration Linux autour de la gestion interactive d'utilisateurs dans un environnement NFS/NIS.

\section{Contexte}

L'environnement repose sur deux distributions Debian sous WSL2 : un serveur NFS/NIS et un client NFS/NIS.

\section{Architecture}

Le serveur centralise les comptes avec NIS et exporte les repertoires personnels avec NFS.

\section{Script interactif}

Cette section decrira la logique du script Bash de gestion des utilisateurs.

\section{Tests}

Cette section presentera les tests realises cote serveur et cote client.

\section{Limites}

Cette section indiquera les limites observees sous WSL.

\section{Conclusion}

\end{document}
